Vision-Language Model (VLM) แบบเปิดสำหรับ OCR ภาษาไทยอ่านหน้าเอกสารได้ถูกต้องเกือบทั้งหน้าแล้ว แต่ยังมีคำถามที่สำคัญในทางปฏิบัติอยู่ข้อหนึ่ง คือ สระและวรรณยุกต์ของไทยมีขนาดเล็ก และเครื่องหมายเพียงตัวเดียวก็เปลี่ยนความหมายของคำได้ งานนี้จึงถามว่าเครื่องหมายเหล่านี้หายไปที่ใด ไม่ได้ถามเพียงว่าหายไปเท่าไร เราเสนอคะแนนที่วัดเครื่องหมายโดยไม่หักคะแนนเรื่องลำดับการอ่าน หักคะแนนข้อความที่เขียนเกิน และไม่ให้คะแนนกับข้อความที่ซ้ำจากการวนซ้ำ บนชุด calibration ของ ThaiOCRBench โมเดล Typhoon OCR 1.5 (2B) เขียนวรรณยุกต์ผิดเพียง 0.65\% ในตำแหน่งที่อ่านพยัญชนะฐานถูกต้อง และเมื่อให้ข้อความก่อนหน้าที่ถูกต้อง โมเดลเลือกวรรณยุกต์ถูกต้อง 98.8\% แม้ไม่เห็นภาพ เครื่องหมายที่หายไปส่วนใหญ่ไม่ได้ถูกอ่านเลย กล่าวคือ บรรทัดหายไปทั้งบรรทัด ข้อความขาดเป็นช่วง หรือ output ถูกการวนซ้ำกินไป ในบรรดาวิธี decode ที่ไม่ต้อง train เพิ่ม การหยุดเมื่อตรวจพบการวนซ้ำไม่เปลี่ยน output ที่ไม่วนซ้ำเลยและเพิ่มคะแนน mark F1 ขึ้น 0.8 จุดสำหรับ Typhoon และ 10.8 จุดสำหรับโมเดลไทยขนาด 0.9B ส่วน repetition penalty ทำให้ F1 ของหน้าที่ไม่วนซ้ำลดลง และการหักความรู้ภาษาออกทำให้วรรณยุกต์ที่ถูกอยู่แล้วหายไป \todo{ผลของโมเดลอื่นและการยืนยันบนชุดข้อมูลที่เก็บไว้ (locked split) ยังไม่เสร็จ} ผลทั้งหมดเป็นผลเบื้องต้น ใช้ benchmark เดียว ชุดข้อมูล calibration และกฎบางข้อกำหนดหลังจากเห็นข้อมูลแล้ว
